%!TEX root = forallxyyc.tex
\part{Interpretações}
\label{ch.semantics}
\addtocontents{toc}{\protect\mbox{}\protect\hrulefill\par}


\chapter{Extensionalidade}\label{s:Interpretations}

Lembre-se de que a LVF é uma lógica verofuncional. Seus conectivos são todos verofuncionais, e \emph{tudo} o que podemos fazer com a LVF é associar sentenças a determinados valores de verdade. Podemos fazer isso \emph{diretamente}. Por exemplo, poderíamos estipular que a sentença da LVF \textquotedblleft $P$\textquotedblright{} deve ser verdadeira. Alternativamente, podemos fazer isso \emph{indiretamente}, oferecendo uma chave de simbolização, por exemplo:
	\begin{ekey}
\item[P] Big Ben está em Londres.
	\end{ekey}
Mas lembre-se de \cref{s:TruthFunctionality} que este é \emph{apenas} um meio de especificar o valor de verdade de \textquotedblleft $P$\textquotedblright{}; a afirmação da chave de simbolização equivale a algo como a seguinte estipulação:
	\begin{itemize}
\item A sentença da LVF \textquotedblleft $P$\textquotedblright{} é verdadeira \ifeff{} Big Ben está em Londres.
	\end{itemize}
E enfatizamos em \cref{s:TruthFunctionality} que a LVF não consegue lidar com diferenças de significado que vão além de meras diferenças de valor de verdade.

\section{Simbolizar versus traduzir}

A LPO tem algumas limitações semelhantes. Ela vai além de meros valores de verdade, pois nos permite decompor sentenças em termos, predicados e quantificadores. Isso nos permite considerar o que é \emph{verdadeiro de} um objeto particular, ou de alguns ou de todos os objetos. \emph{Mas é só isso}.




Para desenvolver um pouco essa ideia, considere esta chave de simbolização:
	\begin{ekey}
\item[\atom{C}{x}] \gap{x} ensina Lógica III em Calgary
	\end{ekey}
Essa estipulação não transporta o \emph{significado} do predicado em inglês para o nosso predicado da LPO. Estamos simplesmente estipulando algo como isto:
	\begin{itemize}
\item \textquotedblleft $\atom{C}{x}$\textquotedblright{} e \textquotedblleft \gap{x} ensina Lógica III em Calgary\textquotedblright{} devem ser \emph{verdadeiras de} exatamente as mesmas coisas.
	\end{itemize}
Assim, em particular:
	\begin{itemize}
\item \textquotedblleft $\atom{C}{x}$\textquotedblright{} deve ser verdadeira exatamente das coisas que ensinam Lógica III em Calgary (quaisquer que sejam essas coisas).
	\end{itemize}
Esta é uma maneira indireta de estipular de quais coisas um predicado é verdadeiro.

Alternativamente, podemos estipular diretamente as extensões dos predicados. Por exemplo, podemos estipular que \textquotedblleft $\atom{C}{x}$\textquotedblright{} deve ser verdadeiro de Richard Zach, e somente de Richard Zach. Acontece que essa estipulação direta teria o mesmo efeito que a estipulação indireta, pois Richard, e somente Richard, ensina Lógica~III em Calgary. Observe, contudo, que os predicados em inglês \textquotedblleft \blank\ é Richard Zach\textquotedblright{} e \textquotedblleft \blank\ ensina Lógica III em Calgary\textquotedblright{} têm significados muito diferentes!

O ponto é que a LPO não tem recursos para lidar com nuances de significado. Quando interpretamos a LPO, tudo o que consideramos é de que coisas os predicados são verdadeiros, independentemente de especificarmos essas coisas direta ou indiretamente. As coisas de que um predicado é verdadeiro são conhecidas como a \define{extens\~ao} desse predicado. Dizemos que a LPO é uma \define{linguagem extensional} porque não representa diferenças de significado entre predicados que têm a mesma extensão.

É por isso que falamos em \emph{simbolizar} sentenças em português na LPO. Se traduzir tiver de preservar o significado, fica claro que não estamos \emph{traduzindo} o português para a LPO.


\section{Extensões}\label{sec:extensions}
Podemos estipular diretamente de que coisas os predicados devem ser verdadeiros. E nossas estipulações podem ser tão arbitrárias quanto quisermos. Por exemplo, poderíamos estipular que \textquotedblleft $\atom{H}{x}$\textquotedblright{} deve ser verdadeiro, e somente verdadeiro, dos seguintes objetos:
	\begin{itemize}
		\item Aristotle
		\item the number $\pi$
\item toda tecla correspondente à nota dó aguda em todo piano já fabricado
	\end{itemize}
Munidos dessa interpretação de \textquotedblleft $\atom{H}{x}$\textquotedblright{}, suponha que acrescentemos à nossa chave de simbolização:
	\begin{ekey}
		\item[a] Aristotle
		\item[h] Hypatia
		\item[p] the number $\pi$
	\end{ekey}
Então \textquotedblleft $\atom{H}{a}$\textquotedblright{} e \textquotedblleft $\atom{H}{p}$\textquotedblright{} serão ambas verdadeiras, nessa interpretação, mas \textquotedblleft $\atom{H}{h}$\textquotedblright{} será falsa, pois Hipátia não estava entre os objetos estipulados.

Esse processo de estipulação explícita é às vezes descrito como estipular a \emph{extensão} de um predicado. Observe que, na estipulação que acabamos de fornecer, os objetos listados não têm nada de particularmente comum. Isso não importa. A lógica não se preocupa com o que nós, humanos (em um determinado momento), pensamos que \textquotedblleft naturalmente combina\textquotedblright{}; para a lógica, todos os objetos estão em pé de igualdade.

Qualquer coleção bem definida de objetos é uma extensão potencial de um predicado de um lugar. O exemplo acima mostra uma maneira de estipular a extensão de \textquotedblleft $\atom{H}{x}$\textquotedblright{} por \emph{enumeração}, isto é, simplesmente listamos os objetos que pertencem à extensão de~\textquotedblleft $\atom{H}{x}$\textquotedblright{}. Também podemos estipular a extensão, como já vimos, fornecendo um predicado em português, como \textquotedblleft \gap{x} ensina Lógica~III em Calgary\textquotedblright{} ou \textquotedblleft \gap{x} é um inteiro par entre $3$ e $9$\textquotedblright{}. O segundo especificaria uma extensão constituída por, e somente por, $4$, $6$ e~$8$.

Observe que alguns predicados do português, como \textquotedblleft \gap{x} é um quadrado redondo\textquotedblright{}, não são verdadeiros de coisa alguma. Nesse caso, dizemos que a extensão do predicado é \emph{vazia}. Permitimos extensões vazias e podemos estipular que a extensão de \textquotedblleft $\atom{H}{x}$\textquotedblright{} seja vazia simplesmente não listando nenhum membro. (Pode ser estranho considerar coleções de nenhuma coisa, mas a lógica às vezes é estranha desse modo.)


\section{Predicados de muitos lugares}
Tudo isso é bastante fácil de entender no caso de predicados de um lugar, mas fica mais complicado quando lidamos com predicados de dois lugares. Considere uma chave de simbolização como:
	\begin{ekey}
\item[\atom{L}{x,y}] \gap{x} ama \gap{y}
	\end{ekey}
Dado o que dissemos acima, essa chave de simbolização deve ser lida como dizendo:
	\begin{itemize}
\item \textquotedblleft $\atom{L}{x,y}$\textquotedblright{} e \textquotedblleft \gap{x} ama \gap{y}\textquotedblright{} devem ser verdadeiros exatamente das mesmas coisas.
	\end{itemize}
Assim, em particular:
	\begin{itemize}
\item \textquotedblleft $\atom{L}{x,y}$\textquotedblright{} deve ser verdadeira de $x$ e $y$ (nessa ordem) \ifeff{} $x$ ama $y$.
	\end{itemize}
É importante insistirmos na ordem aqui, pois o amor — notoriamente — nem sempre é correspondido. (Observe que \textquotedblleft $x$\textquotedblright{} e \textquotedblleft $y$\textquotedblright{} à direita são símbolos do português aumentado e estão sendo \emph{usados}. Em contraste, \textquotedblleft $x$\textquotedblright{} e \textquotedblleft $y$\textquotedblright{} em \textquotedblleft $\atom{L}{x,y}$\textquotedblright{} são símbolos da LPO e estão sendo \emph{mencionados}.)

Essa é uma estipulação indireta. E quanto a uma estipulação direta? Isso também é complicado. Se simplesmente listarmos os objetos que caem sob \textquotedblleft $\atom{L}{x,y}$\textquotedblright{}, não saberemos se são o amante ou o amado (ou ambos). Temos de encontrar uma maneira de incluir a ordem em nossa estipulação explícita.

Para isso, podemos especificar que predicados de dois lugares são verdadeiros de \emph{pares} de objetos, nos quais a ordem do par é importante. Assim, poderíamos estipular que \textquotedblleft $\atom{B}{x,y}$\textquotedblright{} deve ser verdadeiro, e somente verdadeiro, dos seguintes pares de objetos:
	\begin{itemize}
		\item \ntuple{Lenin, Marx}
		\item \ntuple{de Beauvoir, Sartre}
		\item \ntuple{Sartre, de Beauvoir}
	\end{itemize}
Aqui, os colchetes angulares nos mantêm informados sobre a ordem. Suponha que agora acrescentemos as seguintes estipulações:
	\begin{ekey}
		\item[l] Lenin
		\item[m] Marx
		\item[b] de Beauvoir
		\item[r] Sartre
	\end{ekey}
Então \textquotedblleft $\atom{B}{l,m}$\textquotedblright{} será verdadeira, pois \ntuple{Lenin, Marx} está em nossa lista explícita, mas \textquotedblleft $\atom{B}{m,l}$\textquotedblright{} será falsa, pois \ntuple{Marx, Lenin} não está em nossa lista. Entretanto, tanto \textquotedblleft $\atom{B}{b,r}$\textquotedblright{} quanto \textquotedblleft $\atom{B}{r,b}$\textquotedblright{} serão verdadeiras, pois tanto \ntuple{de Beauvoir, Sartre} quanto \ntuple{Sartre, de Beauvoir} estão em nossa lista explícita.

Para tornar essas ideias mais precisas, precisaríamos desenvolver alguma \emph{teoria elementar dos conjuntos}. A teoria dos conjuntos tem um aparato formal que nos permite lidar com extensões, pares ordenados e assim por diante. No entanto, a teoria dos conjuntos não é abordada neste livro. Portanto, deixarei essas ideias em um nível impreciso. Ainda assim, a ideia geral deve estar clara.


\section{Semântica da identidade}
A identidade é um predicado especial da LPO. Nós a escrevemos de maneira um pouco diferente de outros predicados de dois lugares: \textquotedblleft $x=y$\textquotedblright{} em vez de \textquotedblleft $\atom{I}{x,y}$\textquotedblright{} (por exemplo). Mais importante, porém, é que sua interpretação é fixada de uma vez por todas.

Se dois nomes se referem ao mesmo objeto, então trocar um nome pelo outro não mudará o valor de verdade de nenhuma sentença. Assim, em particular, se \textquotedblleft $a$\textquotedblright{} e \textquotedblleft $b$\textquotedblright{} nomeiam o mesmo objeto, todas as seguintes sentenças serão verdadeiras:\label{model.nonidentity}
	\begin{align*}
	 	\atom{A}{a} &\eiff \atom{A}{b} \\
	 	\atom{B}{a} &\eiff \atom{B}{b}\\
		\atom{R}{a,a} &\eiff \atom{R}{b,b}\\
		\atom{R}{a,a} & \eiff \atom{R}{a,b}\\
		\atom{R}{c,a} &\eiff \atom{R}{c,b}\\
		\forall x\, \atom{R}{x,a} &\eiff \forall x\, \atom{R}{x,b}
	\end{align*}
Alguns filósofos acreditaram no inverso dessa afirmação. Isto é, acreditaram que, quando exatamente as mesmas sentenças (que não contenham \textquotedblleft $=$\textquotedblright{}) são verdadeiras de $x$ e~$y$, então $x$ e $y$ são exatamente o mesmo objeto. Essa é uma afirmação filosófica altamente controversa — às vezes chamada de \emph{identidade dos indiscerníveis} — e nossa lógica não a aceitará; permitimos que exatamente as mesmas coisas sejam verdadeiras de dois objetos \emph{distintos}.

Para deixar isso claro, considere a seguinte interpretação:
	\begin{ekey}
\item[\text{domínio}] P.~D.\ Magnus, Tim Button
		\item[a] P.~D.\ Magnus
		\item[b] Tim Button
\item[\atom{\metav{R}}{x_1,\dots, x_n}] Para todo predicado primitivo \metav{R} que decidamos considerar, esse predicado é verdadeiro de \emph{nada}.


	\end{ekey}
Suponha que \textquotedblleft $A$\textquotedblright{} seja um predicado de um lugar; então \textquotedblleft $\atom{A}{a}$\textquotedblright{} é falsa e \textquotedblleft $\atom{A}{b}$\textquotedblright{} é falsa, de modo que \textquotedblleft $\atom{A}{a} \eiff \atom{A}{b}$\textquotedblright{} é verdadeira. Da mesma forma, se \textquotedblleft $R$\textquotedblright{} é um predicado de dois lugares, então \textquotedblleft $\atom{R}{a,a}$\textquotedblright{} é falsa e \textquotedblleft $\atom{R}{a,b}$\textquotedblright{} é falsa, de modo que \textquotedblleft $\atom{R}{a,a} \eiff \atom{R}{a,b}$\textquotedblright{} é verdadeira. E assim por diante: toda sentença atômica que não envolva \textquotedblleft $=$\textquotedblright{} é falsa, portanto todo bicondicional que ligue tais sentenças é verdadeiro. Apesar de tudo, Tim Button e P.~D.\ Magnus são duas pessoas distintas, e não uma única e mesma pessoa!

\section{Interpretações}
Definimos uma \define{valora\c{c}\~ao} na LVF como qualquer atribuição de verdade e falsidade a letras sentenciais. Na LPO, vamos definir uma \define{interpreta\c{c}\~ao} como consistindo de quatro coisas:\footnote{Interpretações também são frequentemente chamadas de \textquotedblleft modelos\textquotedblright{} ou \textquotedblleft estruturas\textquotedblright{}.}



	\begin{compactlist}
\item a especificação de um domínio não vazio;
\item a cada letra sentencial que decidamos considerar é atribuído um valor de verdade;
\item a cada nome que decidamos considerar é atribuído exatamente um objeto dentro do domínio (seu \define{referente});
\item para cada predicado que decidamos considerar (exceto \textquotedblleft $=$\textquotedblright{}), uma especificação de de que coisas (e em que ordem) o predicado é verdadeiro (sua \define{extens\~ao}).



	\end{compactlist}
Para predicados de um lugar, as extensões são coleções de objetos do domínio dos quais o predicado é verdadeiro; para predicados de dois lugares, são pares ordenados de objetos, etc.



Não precisamos especificar nada para~\textquotedblleft $=$\textquotedblright{}, pois ele tem um significado \emph{fixo}, a saber, o da identidade. Tudo é idêntico a si mesmo, e somente a si mesmo.



As chaves de simbolização que consideramos em \cref{ch.FOL} consequentemente nos dão uma maneira muito conveniente de apresentar uma interpretação. Continuaremos a usá-las neste capítulo. Seguindo a discussão de \cref{sec:extensions}, agora também permitimos extensões especificadas por enumerações no lado direito, por exemplo,
\begin{ekey}
\item[\text{domínio}] filósofos, números
	\item[\atom{H}{x}] Aristotle, Hypatia, $\pi$
\end{ekey}
é uma maneira perfeitamente boa de especificar uma interpretação, assim como
\begin{ekey}
\item[\text{domínio}] $0$, $1$, $2
	\item[\atom{L}{x,y}] \ntuple{$0$,$1$}, \ntuple{$0$, $2$}, \ntuple{$1$, $2$}
\end{ekey}
Poderíamos ter especificado a mesma extensão (neste domínio particular) fornecendo o predicado em português \textquotedblleft \gap{x} é menor que \gap{y}\textquotedblright{}.


\newglossaryentry{interpretation}{
  name = {interpretação},
  description = {Uma especificação de um \gls{domain} juntamente com os objetos que os \glspl{name} designam e dos quais os \glspl{predicate} são verdadeiros}
}

No entanto, às vezes também é conveniente apresentar uma interpretação \emph{diagramaticamente}. Para ilustrar (literalmente): suponha que queiramos considerar apenas um predicado de dois lugares, \textquotedblleft $\atom{R}{x,y}$\textquotedblright{}. Podemos representá-lo simplesmente desenhando uma seta entre dois objetos e estipulando que \textquotedblleft $\atom{R}{x,y}$\textquotedblright{} vale para $x$ e $y$ \ifeff{} há uma seta que vai de $x$ a $y$ em nosso diagrama. Como exemplo, poderíamos oferecer:





\begin{center}
\begin{tikzpicture}[alt={Os números 1, 2, 3 e 4 estão conectados por setas. Há setas de 1 para 2 e 3, de 2 para 3, de 3 para 4 e de 4 para 1.}]


\node (atom1) at (0,2) {$1$};
\node (atom2) at (2,2) {$2$};
\node (atom3) at (2,0) {$3$};
\node (atom4) at (0,0) {$4$};
\draw[->, thick] (atom1)--(atom2);
\draw[->, thick] (atom2)--(atom3);
\draw[->, thick] (atom3)--(atom4);
\draw[->, thick] (atom4)--(atom1);
\draw[->, thick] (atom1) -- (atom3);
\end{tikzpicture}
\bmlDescription{Os números 1, 2, 3 e 4 estão conectados por setas. Há setas de 1 para 2 e 3, de 2 para 3, de 3 para 4 e de 4 para 1.}


\end{center}
Este diagrama poderia ser usado para descrever uma interpretação cujo domínio são os quatro primeiros números inteiros positivos e que interpreta \textquotedblleft $\atom{R}{x,y}$\textquotedblright{} como verdadeiro de, e somente de:
	\begin{center}
		\ntuple{$1$, $2$},
		\ntuple{$2$, $3$},
		\ntuple{$3$, $4$},
		\ntuple{$4$, $1$},
		\ntuple{$1$, $3$}
	\end{center}
De modo equivalente, poderíamos oferecer este diagrama:

\begin{center}
\begin{tikzpicture}[alt={Os números 1, 3 e 4 estão conectados por setas. Há setas de 1 para 3 e para si mesmo, e de 3 para 1, 4 e para si mesmo. O número 2 está incluído, mas não há setas entrando ou saindo dele.}]

\node (atom1) at (0,2) {$1$};
\node (atom2) at (2,2) {$2$};
\node (atom3) at (2,0) {$3$};
\node (atom4) at (0,0) {$4$};
\draw[->, thick] (atom3)--(atom4);
\draw[->, thick] (atom1)+(-0.15,0.15) arc (-330:-30:.3);
\draw[->, thick] (atom3)+(0.15,-0.15) arc (-150:150:.3);
\draw[<->, thick] (atom1) -- (atom3);
\end{tikzpicture}
\bmlDescription{Os números 1, 3 e 4 estão conectados por setas. Há setas de 1 para 3 e para si mesmo, e de 3 para 1, 4 e para si mesmo. O número 2 está incluído, mas não há setas entrando ou saindo dele.}

\end{center}

A interpretação especificada por este diagrama também pode ser dada listando o que está no domínio e na extensão de~\textquotedblleft $\atom{R}{x,y}$\textquotedblright{}:
	\begin{ekey}
\item[\text{domínio}] $1$, $2$, $3$, $4
	\item[R(x,y)] \ntuple{$1$, $3$},
		\ntuple{$3$, $1$},
		\ntuple{$3$, $4$},
		\ntuple{$1$, $1$},
		\ntuple{$3$, $3$}
	\end{ekey}
Se quiséssemos, poderíamos tornar nossos diagramas mais complexos. Por exemplo, poderíamos acrescentar nomes como rótulos para objetos particulares. Da mesma forma, para simbolizar a extensão de um predicado de um lugar, poderíamos simplesmente desenhar um círculo em torno de alguns objetos particulares e estipular que os objetos assim circundados (e somente eles) caem sob o predicado \textquotedblleft $\atom{H}{x}$\textquotedblright{}, digamos. Para especificar múltiplos predicados, poderíamos usar linhas coloridas (ou tracejadas, pontilhadas) para setas e círculos.


\chapter{Verdade na LPO}\label{s:TruthFOL}
Nós o apresentamos às interpretações. Como, entre outras coisas, elas nos dizem de que objetos os predicados são verdadeiros, fornecerão uma explicação da verdade de sentenças atômicas. No entanto, agora precisamos dizer, precisamente, o que é uma sentença arbitrária da LPO ser verdadeira ou falsa em uma interpretação.

Sabemos de \cref{s:FOLSentences} que há três tipos de sentença na LPO:
	\begin{compactlist}
\item sentenças atômicas,
\item sentenças cujo operador lógico principal é um conectivo sentencial,
\item sentenças cujo operador lógico principal é um quantificador.
	\end{compactlist}
Precisamos explicar a verdade para os três tipos de sentença.

Nesta seção, ofereceremos uma explicação completamente geral. No entanto, para tentar manter a explicação compreensível, usaremos, em vários pontos, a seguinte interpretação:
	\begin{ekey}
\item[\text{domínio}] todas as pessoas nascidas antes de 2000 \textsc{d.C.}
		\item[a] Aristotle
		\item[b] Beyonc\'e
\item[\atom{P}{x}] \gap{x} é filósofo
\item[\atom{R}{x,y}] \gap{x} nasceu antes de \gap{y}
	\end{ekey}
Este será nosso \emph{exemplo de referência} no que segue.

\section{Sentenças atômicas}
A verdade de sentenças atômicas deve ser bastante direta. Para letras sentenciais, a interpretação especifica se elas são verdadeiras ou falsas. A sentença \textquotedblleft $\atom{P}{a}$\textquotedblright{} deve ser verdadeira se, e somente se, \textquotedblleft $\atom{P}{x}$\textquotedblright{} for verdadeira de \textquotedblleft $a$\textquotedblright{}. Dada nossa interpretação de referência, isso é verdadeiro \ifeff{} Aristóteles é filósofo. Aristóteles é filósofo. Portanto, a sentença é verdadeira. Da mesma forma, \textquotedblleft $\atom{P}{b}$\textquotedblright{} é falsa em nossa interpretação de referência.

Da mesma forma, nessa interpretação, \textquotedblleft $\atom{R}{a,b}$\textquotedblright{} é verdadeira \ifeff{} o objeto nomeado por \textquotedblleft $a$\textquotedblright{} nasceu antes do objeto nomeado por \textquotedblleft $b$\textquotedblright{}. Ora, Aristóteles nasceu antes de Beyonc\'e. Portanto, \textquotedblleft $\atom{R}{a,b}$\textquotedblright{} é verdadeira. Do mesmo modo, \textquotedblleft $\atom{R}{a,a}$\textquotedblright{} é falsa: Aristóteles não nasceu antes de Aristóteles.

Lidar com sentenças atômicas, então, é muito intuitivo. Quando \metav{R} é um predicado de $n$ lugares e $\metav{a}_1$, $\metav{a}_{2}$, \dots, $\metav{a}_{n}$ são nomes,

	\factoidbox{
		A sentença $\atom{\metav{R}}{\metav{a}_{1},\metav{a}_{2},\dots,\metav{a}_{n}}$ é verdadeira em uma interpretação \textbf{\ifeff}
		$\metav{R}$ é verdadeiro dos objetos nomeados por $\metav{a}_{1}$, $\metav{a}_{2}$, \dots, $\metav{a}_{n}$ (nessa ordem) nessa interpretação.
	}
Lembre-se, porém, de que há um tipo especial de sentença atômica: dois nomes conectados por um sinal de identidade constituem uma sentença atômica. Esse tipo de sentença atômica também é fácil de tratar. Sejam \metav{a} e \metav{b} nomes quaisquer,
	\factoidbox{
		$\metav{a} = \metav{b}$ é verdadeira em uma interpretação \textbf{\ifeff}
		\metav{a} e \metav{b} nomeiam exatamente o mesmo objeto nessa interpretação.
	}
Assim, em nossa interpretação de referência, \textquotedblleft $a = b$\textquotedblright{} é falsa, pois Aristóteles é distinto de Beyonc\'e.


\section{Conectivos sentenciais}
Vimos em \cref{s:FOLSentences} que sentenças da LPO podem ser construídas a partir de sentenças mais simples usando os conectivos verofuncionais familiares da LVF. As regras que governam esses conectivos verofuncionais são \emph{exatamente} as mesmas que eram quando consideramos a LVF. Ei-las:
	\factoidbox{
		\begin{factoiditems}
\item $\metav{A} \eand \metav{B}$ é verdadeira em uma interpretação \textbf{\ifeff} tanto $\metav{A}$ quanto $\metav{B}$ são verdadeiras nessa interpretação.

\item $\metav{A} \eor \metav{B}$ é verdadeira em uma interpretação \textbf{\ifeff} ou $\metav{A}$ ou $\metav{B}$ é verdadeira nessa interpretação.

\item $\enot \metav{A}$ é verdadeira em uma interpretação \textbf{\ifeff} $\metav{A}$ é falsa nessa interpretação.

\item $\metav{A} \eif \metav{B}$ é verdadeira em uma interpretação \textbf{\ifeff} ou $\metav{A}$ é falsa ou $\metav{B}$ é verdadeira nessa interpretação.

\item $\metav{A} \eiff \metav{B}$ é verdadeira em uma
		interpretação \textbf{\ifeff} $\metav{A}$ tem o mesmo valor de verdade
		que $\metav{B}$ nessa interpretação.
		\end{factoiditems}
	}
Isso apresenta exatamente a mesma informação que as tabelas-verdade características dos conectivos; apenas o faz de uma maneira um pouco diferente. Alguns exemplos provavelmente ajudarão a ilustrar a ideia. (Certifique-se de compreendê-los!) Em nossa interpretação de referência:
	\begin{itemize}
\item \textquotedblleft $a = a \eand \atom{P}{a}$\textquotedblright{} é verdadeira.
\item \textquotedblleft $\atom{R}{a,b} \eand \atom{P}{b}$\textquotedblright{} é falsa porque, embora \textquotedblleft $\atom{R}{a,b}$\textquotedblright{} seja verdadeira, \textquotedblleft $\atom{P}{b}$\textquotedblright{} é falsa.
\item \textquotedblleft $a = b \eor \atom{P}{a}$\textquotedblright{} é verdadeira.
\item \textquotedblleft $\enot a = b$\textquotedblright{} é verdadeira.
\item \textquotedblleft $\atom{P}{a} \eand \enot( a= b \eand \atom{R}{a,b})$\textquotedblright{} é verdadeira, porque \textquotedblleft $\atom{P}{a}$\textquotedblright{} é verdadeira e \textquotedblleft $a = b$\textquotedblright{} é falsa.
	\end{itemize}
Certifique-se de que compreende esses exemplos.

\ifHTMLtarget
\section{Quando o operador lógico principal é um quantificador}
\else
\section[Quantificadores]{Quando o operador lógico principal é um quantificador}
\fi
\label{s:MainLogicalOperatorQuantifier}

A inovação empolgante da LPO, porém, é o uso de \emph{quantificadores}, mas expressar as condições de verdade de sentenças quantificadas é um pouco mais trabalhoso do que se poderia esperar inicialmente.

Eis um primeiro pensamento ingênuo. Queremos dizer que \textquotedblleft $\forall x\, \atom{F}{x}$\textquotedblright{} é verdadeira \ifeff{} \textquotedblleft $\atom{F}{x}$\textquotedblright{} é verdadeira de tudo no domínio. Isso não deveria ser muito problemático: nossa interpretação especificará diretamente de que coisas \textquotedblleft $\atom{F}{x}$\textquotedblright{} é verdadeira.

Infelizmente, esse pensamento ingênuo não é geral o bastante. Por exemplo, queremos poder dizer que \textquotedblleft $\forall x \exists y\, \atom{L}{x,y}$\textquotedblright{} é verdadeira se, grosso modo, \textquotedblleft $\exists y\, \atom{L}{x,y}$\textquotedblright{} é verdadeira de tudo no domínio. Mas nossa interpretação não especifica \emph{diretamente} de que coisas \textquotedblleft $\exists y\, \atom{L}{x,y}$\textquotedblright{} é verdadeira. Em vez disso, se isso é ou não verdadeiro de algo deve resultar apenas da interpretação do predicado \textquotedblleft $L$\textquotedblright{}, do domínio e dos significados dos quantificadores.

Eis então um segundo pensamento ingênuo. Poderíamos tentar dizer que \textquotedblleft $\forall x \exists y\, \atom{L}{x,y}$\textquotedblright{} é verdadeira em uma interpretação \ifeff{} $\exists y\, \atom{L}{\metav{a},y}$ é verdadeira para \emph{todo} nome \metav{a} que incluímos em nossa interpretação. Da mesma forma, poderíamos tentar dizer que $\exists y\, \atom{L}{\metav{a},y}$ é verdadeira se, e somente se, $\atom{L}{\metav{a},\metav{b}}$ é verdadeira para \emph{algum} nome \metav{b} que incluímos em nossa interpretação.

Infelizmente, isso também não está correto. Para ver isso, observe que nossa interpretação de referência interpreta apenas \emph{dois} nomes, \textquotedblleft $a$\textquotedblright{} e \textquotedblleft $b$\textquotedblright{}. Mas o domínio — todas as pessoas nascidas antes do ano 2000 \textsc{d.C.} — contém muito mais de duas pessoas. (E não pretendemos corrigir isso nomeando \emph{todas} elas!)

Eis então um terceiro pensamento. (E este pensamento não é ingênuo, mas correto.) Embora não tenhamos nomeado \emph{todos}, cada pessoa \emph{poderia} ter recebido um nome. Portanto, devemos concentrar-nos nessa possibilidade de estender uma interpretação acrescentando um novo nome. Ofereceremos alguns exemplos de como isso poderia funcionar, centrados em nossa interpretação de referência, e então apresentaremos a definição formal.

Em nossa interpretação de referência, \textquotedblleft $\exists x\, \atom{R}{b,x}$\textquotedblright{} deve ser verdadeira. Afinal, no domínio, certamente há alguém que nasceu depois de Beyonc\'e. Lady Gaga é uma dessas pessoas. De fato, se estendêssemos nossa interpretação de referência — temporariamente, é claro — acrescentando o novo nome~\textquotedblleft $c$\textquotedblright{} para referir-se a Lady Gaga, então \textquotedblleft $\atom{R}{b,c}$\textquotedblright{} seria verdadeira nessa interpretação estendida. Isso, certamente, deve ser suficiente para tornar \textquotedblleft $\exists x\, \atom{R}{b,x}$\textquotedblright{} verdadeira na interpretação de referência original.
verdadeira. Afinal, no domínio, há certamente alguém que nasceu
depois de Beyonc\'e. Lady Gaga é uma dessas pessoas. De fato, se
estendêssemos nossa interpretação de referência — temporariamente,
é claro — acrescentando o novo nome~\textquotedblleft $c$\textquotedblright{} para referir-se a Lady Gaga, então
\textquotedblleft $\atom{R}{b,c}$\textquotedblright{} seria verdadeira nessa interpretação estendida. Isso,
certamente, deve ser suficiente para tornar \textquotedblleft $\exists x\, \atom{R}{b,x}$\textquotedblright{} verdadeira na
interpretação original de referência.

Em nossa interpretação de referência, \textquotedblleft $\exists x (\atom{P}{x} \eand \atom{R}{x,a})$\textquotedblright{} também deve ser verdadeira. Afinal, no domínio, certamente há alguém que é filósofo e nasceu antes de Aristóteles. Sócrates é uma dessas pessoas. De fato, se estendêssemos nossa interpretação de referência fazendo com que um novo nome, \textquotedblleft $c$\textquotedblright{}, designasse Sócrates, então \textquotedblleft $\atom{P}{c} \eand \atom{R}{c,a}$\textquotedblright{} seria verdadeira nessa interpretação estendida. Novamente, isso deve ser suficiente para tornar \textquotedblleft $\exists x (\atom{P}{x} \eand \atom{R}{x,a})$\textquotedblright{} verdadeira na interpretação original.

Em nossa interpretação de referência, \textquotedblleft $\forall x \exists y\, \atom{R}{x,y}$\textquotedblright{} deve ser falsa. Afinal, considere a última pessoa nascida no ano de 1999. Não sabemos quem foi, mas, se estendêssemos nossa interpretação de referência fazendo com que um novo nome, \textquotedblleft $d$\textquotedblright{}, designasse essa pessoa, não conseguiríamos encontrar outra pessoa no domínio para designar com algum novo nome adicional, talvez \textquotedblleft $e$\textquotedblright{}, de tal modo que \textquotedblleft $\atom{R}{d,e}$\textquotedblright{} fosse verdadeira. Na verdade, não importa \emph{quem} nomeássemos com \textquotedblleft $e$\textquotedblright{}, \textquotedblleft $\atom{R}{d,e}$\textquotedblright{} seria falsa. Essa observação é certamente suficiente para tornar \textquotedblleft $\exists y\, \atom{R}{d,y}$\textquotedblright{} \emph{falsa} na interpretação estendida, o que, por sua vez, é certamente suficiente para tornar \textquotedblleft $\forall x \exists y\, \atom{R}{x,y}$\textquotedblright{} falsa na interpretação original.

Se você compreendeu esses três exemplos, isso é o que importa. Eles fornecem a base para uma definição formal de verdade de sentenças quantificadas.

Estritamente falando, porém, ainda precisamos \emph{dar} essa definição. O resultado, infelizmente, é um pouco feio e exige algumas definições novas. Prepare-se!

Suponha que \metav{A} seja uma fórmula. Escreveremos
$$\metav{A}(\ldots \metav{x} \ldots \metav{x} \ldots)$$ para indicar
que \metav{A} contém alguma quantidade de ocorrências livres da variável~\metav{x} (geralmente uma ou mais, mas não excluímos o caso em que \metav{x} não ocorra livre de modo algum). Suponha também que \metav{c} seja um nome. Então escreveremos:



$$\metav{A}(\ldots \metav{c} \ldots \metav{c} \ldots)$$
para a fórmula que obtemos substituindo \emph{toda} ocorrência livre de $\metav{x}$ em \metav{A} por~$\metav{c}$. A fórmula resultante é chamada de \define{inst\^ancia de substitui\c{c}\~ao} de $\forall \metav{x}\metav{A}$ e $\exists\metav{x}\metav{A}$. Além disso, $\metav{c}$ é chamado de \define{nome instanciador}. (Se \metav{A} não contiver nenhuma ocorrência livre de~$\metav{x}$, então a instância de substituição é idêntica a~\metav{A}. Assim, por exemplo,






	$$\exists x (\atom{R}{e,x} \eiff \atom{F}{x})$$
é uma instância de substituição de
	$$\forall y \exists x (\atom{R}{y,x} \eiff \atom{F}{x})$$
com o nome instanciador \textquotedblleft $e$\textquotedblright{} e a variável instanciada \textquotedblleft $y$\textquotedblright{}.

\newglossaryentry{substitution instance}{
  name = instância de substituição,
  description = {O resultado de substituir toda ocorrência livre de uma \gls{variable} em uma \gls{formula} por um \gls{name}}
}

Uma breve observação: ao substituir nomes por variáveis, permitimos apenas a substituição de variáveis \emph{livres}. Isso é necessário, pois permitimos fórmulas como \textquotedblleft $\exists x(\forall x\,\atom{F}{x} \eif \atom{G}{x})$\textquotedblright{}, nas quais diferentes quantificadores ligam as mesmas variáveis.



Consideramos apenas \textquotedblleft $\forall x\,\atom{F}{x} \eif \atom{G}{e}$\textquotedblright{} como instância de substituição dessa fórmula, e não \textquotedblleft $\atom{F}{e} \eif \atom{G}{e}$\textquotedblright{}, porque o \textquotedblleft $x$\textquotedblright{} em \textquotedblleft $\atom{F}{x}$\textquotedblright{} não é livre — é ligado por \textquotedblleft $\forall x$\textquotedblright{}. (Talvez você queira revisar \cref{s:FreeVars} nesse contexto.)





Nossa interpretação incluirá uma especificação de quais nomes correspondem a quais objetos do domínio. Tome qualquer objeto do domínio, digamos $d$, e um nome $\metav{c}$ que ainda não esteja atribuído pela interpretação. Se nossa interpretação é $\mathbf{I}$, então podemos considerar a interpretação $\mathbf{I}[\text{d}/\metav{c}]$, que é igual à $\mathbf{I}$, exceto que também atribui o objeto~$d$ ao nome $\metav{c}$. Então podemos dizer que $d$ \define{satisfaz} a fórmula







$\metav{A}(\dots\metav{x}\dots\metav{x}\dots)$ na
interpretação~$\mathbf{I}$ se, e somente se,
$\metav{A}(\dots\metav{c}\dots\metav{c}\dots)$ é verdadeira em
$\mathbf{I}[\text{d}/\metav{c}]$. (Supomos que o nome $\metav{c}$ não
ocorra já em
$\metav{A}(\dots\metav{x}\dots\metav{x}\dots)$. Se $d$ satisfaz
$\metav{A}(\dots\metav{x}\dots\metav{x}\dots)$, também dizemos que
$\metav{A}(\dots\metav{x}\dots\metav{x}\dots)$ é \emph{verdadeira de}~$d$.

\factoidbox{
	\begin{factoiditems}
\item A interpretação $\mathbf{I}[\text{d}/\metav{c}]$ é igual à interpretação $\mathbf{I}$, exceto por também atribuir o objeto~$d$ ao nome~$\metav{c}$.



\item Um objeto $d$ \define{satisfaz}
$\metav{A}(\dots\metav{x}\dots\metav{x}\dots)$ em
uma interpretação~$\mathbf{I}$ \textbf{\ifeff}
$\metav{A}(\dots\metav{c}\dots\metav{c}\dots)$ é verdadeira em
$\mathbf{I}[\text{d}/\metav{c}]$ (quando o nome $\metav{c}$ não
ocorre já em $\metav{A}(\dots\metav{x}\dots\metav{x}\dots)$).
	\end{factoiditems}
}

\newglossaryentry{satisfaction}
{
  name=satisfação de uma fórmula,
  description={A relação entre um objeto $d$ no \gls{domain} de uma
    \gls{interpretation}~$\mathbf{I}$ e uma
    \gls{formula}~$\metav{A}(\metav{x})$ que ocorre \ifeff{}
    $\metav{A}(\metav{c})$ é verdadeira na interpretação modificada
    $\mathbf{I}[\text{d}/\metav{c}]$, na qual $\metav{c}$ nomeia~$d$}
}

Assim, por exemplo, Sócrates satisfaz a fórmula~$\atom{P}{x}$, pois $\atom{P}{c}$ é verdadeira na interpretação $\mathbf{I}[\text{Sócrates}/c]$, isto é, na interpretação:
\begin{ekey}
\item[\text{domínio}] todas as pessoas nascidas antes de 2000 \textsc{d.C.}
	\item[a] Aristotle
	\item[b] Beyonc\'e
	\item[c] Socrates
\item[\atom{P}{x}] \gap{x} é filósofo
\item[\atom{R}{x,y}] \gap{x} nasceu antes de \gap{y}
\end{ekey}

Munidos dessa notação, a ideia geral é a seguinte. A sentença
$\forall \metav{x}\metav{A}(\ldots \metav{x} \ldots \metav{x} \ldots)$
será verdadeira em $\mathbf{I}$ \ifeff{}, para qualquer objeto~$d$ no domínio,
$\metav{A}(\ldots \metav{c} \ldots \metav{c}\ldots)$ é verdadeira em
$\mathbf{I}[\text{d}/\metav{c}]$, isto é, não importa qual objeto (no
domínio) nomeamos pelo novo nome~$\metav{c}$. Em outras palavras,
$\forall \metav{x} \metav{A}(\ldots \metav{x} \ldots \metav{x}
\ldots)$ é verdadeira \ifeff{} todo objeto no domínio satisfaz
$\metav{A}(\ldots \metav{x} \ldots \metav{x} \ldots)$. Da mesma forma, a
sentença $\exists \metav{x}\metav{A}(\ldots \metav{x} \ldots \metav{x}
\ldots)$ será verdadeira em $\mathbf{I}$ \ifeff{} existe \emph{algum}
objeto que satisfaz $\metav{A}(\ldots \metav{x} \ldots \metav{x}
\ldots)$, isto é, $\metav{A}(\ldots \metav{c} \ldots \metav{c} \ldots)$
é verdadeira em $\mathbf{I}[\text{d}/\metav{c}]$ para algum objeto~$d$ e um
novo nome~$\metav{c}$. \factoidbox{
		\begin{factoiditems}
\item $\forall \metav{x}\metav{A}(\ldots \metav{x}\ldots\metav{x}\ldots)$ é verdadeira em uma interpretação \textbf{\ifeff}
todo objeto no domínio satisfaz $\metav{A}(\ldots
		$\metav{x} \ldots \metav{x}\ldots)$.

\item $\exists \metav{x}\metav{A}(\ldots \metav{x}\ldots\metav{x}\ldots)$ é verdadeira em uma interpretação \textbf{\ifeff} pelo menos um objeto no domínio satisfaz
		$\metav{A}(\ldots \metav{x}\ldots\metav{x}\ldots)$.
		\end{factoiditems}
	}
Para ser claro: tudo isso está apenas formalizando (de maneira muito pedante) a ideia intuitiva expressa acima. O resultado é um pouco feio e a definição final pode parecer um tanto opaca. Esperamos, porém, que o \emph{espírito} da ideia esteja claro.

\practiceproblems
\solutions
\problempart
\label{pr.TorF1}
Considere a seguinte interpretação:
	\begin{itemize}
\item O domínio compreende apenas Corwin e Benedict
\item \textquotedblleft $\atom{A}{x}$\textquotedblright{} deve ser verdadeira de Corwin e Benedict
\item \textquotedblleft $\atom{B}{x}$\textquotedblright{} deve ser verdadeira apenas de Benedict
\item \textquotedblleft $\atom{N}{x}$\textquotedblright{} deve ser verdadeira de ninguém
\item \textquotedblleft $c$\textquotedblright{} deve referir-se a Corwin
	\end{itemize}
Determine se cada uma das seguintes sentenças é verdadeira ou falsa nessa interpretação:
\begin{compactlist}
\item $\atom{B}{c} $
\item $\atom{A}{c}  \eiff \enot \atom{N}{c}$
\item $\atom{N}{c}  \eif (\atom{A}{c} \eor \atom{B}{c})$
\item $\forall x\, \atom{A}{x}$
\item $\forall x \enot \atom{B}{x}$
\item $\exists x(\atom{A}{x} \eand \atom{B}{x})$
\item $\exists x(\atom{A}{x} \eif \atom{N}{x})$
\item $\forall x(\atom{N}{x} \eor \enot \atom{N}{x})$
\item $\exists x\, \atom{B}{x} \eif \forall x\, \atom{A}{x}$
\end{compactlist}

\problempart
\label{pr.TorF2}
Considere a seguinte interpretação:
	\begin{itemize}
\item O domínio compreende apenas Lemmy, Courtney e Eddy
\item \textquotedblleft $\atom{G}{x}$\textquotedblright{} deve ser verdadeira de Lemmy, Courtney e Eddy.
\item \textquotedblleft $\atom{H}{x}$\textquotedblright{} deve ser verdadeira de, e somente de, Courtney
\item \textquotedblleft $\atom{M}{x}$\textquotedblright{} deve ser verdadeira de, e somente de, Lemmy e Eddy
\item \textquotedblleft $c$\textquotedblright{} deve referir-se a Courtney
\item \textquotedblleft $e$\textquotedblright{} deve referir-se a Eddy
	\end{itemize}
Determine se cada uma das seguintes sentenças é verdadeira ou falsa nessa interpretação:
\begin{compactlist}
\item $\atom{H}{c} $
\item $\atom{H}{e} $
\item $\atom{M}{c}  \eor \atom{M}{e}$
\item $\atom{G}{c}  \eor \enot \atom{G}{c}$
\item $\atom{M}{c}  \eif \atom{G}{c}$
\item $\exists x\, \atom{H}{x}$
\item $\forall x\, \atom{H}{x}$
\item $\exists x\, \enot \atom{M}{x}$
\item $\exists x(\atom{H}{x} \eand \atom{G}{x})$
\item $\exists x(\atom{M}{x} \eand \atom{G}{x})$
\item $\forall x(\atom{H}{x} \eor \atom{M}{x})$
\item $\exists x\, \atom{H}{x} \eand \exists x\, \atom{M}{x}$
\item $\forall x(\atom{H}{x} \eiff \enot \atom{M}{x})$
\item $\exists x\, \atom{G}{x} \eand \exists x \enot \atom{G}{x}$
\item $\forall x\exists y(\atom{G}{x} \eand \atom{H}{y})$
\end{compactlist}

\problempart
\label{pr.TorF3}
Seguindo as convenções diagramáticas introduzidas ao final de \cref{s:Interpretations}, considere a seguinte interpretação:
\begin{center}
\begin{tikzpicture}[alt={Os números 1, 2, 3, 4 e 5 estão conectados por setas. Há setas de 1 para 2, 3, 4, 5 e para si mesmo; de 2 para 5; de 5 para si mesmo; e de 4 para 1. Não há nenhuma seta saindo de 3.}]
\node (atom1) at (0,2) {$1$};
\node (atom2) at (2,2) {$2$};
\node (atom4) at (0,0) {$3$};
\node (atom5) at (2,0) {$4$};
\node (atom6) at (4,0) {$5$};
\draw[->, thick] (atom1)+(-0.15,0.15) arc (-330:-30:.3);
\draw[->, thick] (atom6)+(0.15,-0.15) arc (-150:150:.3);
\draw[->, thick] (atom1) -- (atom2);
\draw[->, thick] (atom1) -- (atom4);
\draw[<->, thick] (atom1) -- (atom5);
\draw[->, thick] (atom1) -- (atom6);
\draw[->, thick] (atom2) -- (atom6);
\end{tikzpicture}
\bmlDescription{Os números 1, 2, 3, 4 e 5 estão conectados por setas. Há setas de 1 para 2, 3, 4, 5 e para si mesmo; de 2 para 5; de 5 para si mesmo; e de 4 para 1. Não há nenhuma seta saindo de 3.}
\end{center}
Determine se cada uma das seguintes sentenças é verdadeira ou falsa nessa interpretação:
\begin{compactlist}
\item $\exists x\, \atom{R}{x,x}$
\item $\forall x\, \atom{R}{x,x}$
\item $\exists x \forall y\, \atom{R}{x,y}$
\item $\exists x \forall y\, \atom{R}{y,x}$
\item $\forall x \forall y \forall z ((\atom{R}{x,y} \eand \atom{R}{y,z}) \eif \atom{R}{x,z})$
\item $\forall x \forall y \forall z ((\atom{R}{x,y} \eand \atom{R}{x,z}) \eif \atom{R}{y,z})$
\item $\exists x \forall y\, \enot \atom{R}{x,y}$
\item $\forall x(\exists y\, \atom{R}{x,y} \eif \exists y\, \atom{R}{y,x})$
\item $\exists x \exists y (\enot x = y \eand \atom{R}{x,y} \eand \atom{R}{y,x})$
\item $\exists x \forall y(\atom{R}{x,y} \eiff x = y)$
\item $\exists x \forall y(\atom{R}{y,x} \eiff x = y)$
\item $\exists x \exists y(\enot x = y \eand \atom{R}{x,y} \eand \forall z(\atom{R}{z,x} \eiff y = z))$
\end{compactlist}


\chapter{Semantic concepts}

\section{Validities, entailment, etc.}

Defining truth in FOL was quite fiddly. But now that we are done, we can define various other central logical notions. These definitions will look very similar to those for TFL, from \cref{s:SemanticConcepts}. However, remember that they concern \emph{interpretations}, rather than valuations.

We will use the symbol `$\entails$' for FOL much as we did for TFL. So
the sentences $\metav{A}_1, \metav{A}_2, \ldots, \metav{A}_n$
\define{entail} the sentence $\metav{C}$,
	$$\metav{A}_1, \metav{A}_2, \ldots, \metav{A}_n \entails\metav{C},$$
\ifeff{} there is no interpretation in which all of $\metav{A}_1$, $\metav{A}_2$, \dots, $\metav{A}_n$ are true and in which \metav{C} is false. Derivatively,
	$$\entails\metav{A}$$
means that \metav{A} is true in every interpretation.

The other logical notions also have corresponding definitions in FOL:

\factoidbox{
	\begin{factoiditems}
		\item An FOL sentence $\metav{A}$ is a \define{validity} \ifeff{}
		$\metav{A}$ is true in every interpretation; i.e.,
		$\entails\metav{A}$.
\newglossaryentry{validity}
{
name=validity,
description={A \gls{sentence of FOL} that is true in every \gls{interpretation}}
}

		\item $\metav{A}$ is a \define{contradiction} \ifeff{} $\metav{A}$
		is false in every interpretation; i.e., $\entails\enot\metav{A}$.
\newglossaryentry{contradiction of FOL}
{
  name=contradiction (of FOL),
  text=contradiction,
description={A \gls{sentence of FOL} that is false in every \gls{interpretation}}
}
  
		\item $\metav{A}_1, \metav{A}_2, \ldots, \metav{A}_n \therefore
		\metav{C}$ is \define{valid in FOL} \ifeff{} there is no
		interpretation in which all of the premises are true and the
		conclusion is false; i.e., $\metav{A}_1,\metav{A}_2,\ldots,
		\metav{A}_n \entails\metav{C}$. It is \define{invalid in FOL}
		otherwise.
\newglossaryentry{valid in FOL}
{
  name=validity of arguments (in FOL),
  text = valid,
description={A property held by arguments; an argument is valid if and only if no \gls{interpretation} makes all premises true and the conclusion false}
}

		\item Two FOL sentences \metav{A} and \metav{B} are
		\define{equivalent} \ifeff{} they are true in exactly the same
		interpretations as each other; i.e., both
		$\metav{A}\entails\metav{B}$ and $\metav{B}\entails\metav{A}$.

\newglossaryentry{equivalent in FOL}
{
  name=equivalence (in FOL),
  text = equivalent,
description={A property held by pairs of \glspl{sentence of FOL} if and only if the sentences have the same truth value in every \gls{interpretation}}
}

		\item The FOL sentences $\metav{A}_1$, $\metav{A}_2$, \dots,
		$\metav{A}_n$ are \define{jointly satisfiable} \ifeff{} some
		interpretation makes all of them true. They are \define{jointly
		unsatisfiable} \ifeff{} there is no such interpretation.
\newglossaryentry{satisfiable in FOL}
{
  name=satisfiability (in FOL),
  text=jointly satisfiable,
  description={A property held by \glspl{sentence of FOL} if and only
  if some \gls{interpretation} makes all the sentences true together}
}
	\end{factoiditems}
}

Note that we use the standard term `validity' for sentences that are
true in every interpretation. Validities are to FOL what tautologies
are to TFL.

\section{Expressibility}

The concept of an object satisfying a formula with one free variable
introduced in \cref{s:MainLogicalOperatorQuantifier} can also be
extended to formulas with more than one free variable. If we have a
formula $\metav{A}(\metav{x},\metav{y})$ with two free variables
$\metav{x}$ and $\metav{y}$, then we can say that a pair of objects
$\langle a, b\rangle$ satisfies $\metav{A}(\metav{x},\metav{y})$
\ifeff{} $\metav{A}(\metav{c},\metav{d})$ is true in the
interpretation extended by two names $\metav{c}$ and $\metav{d}$,
where $\metav{c}$ names~$a$ and $\metav{d}$ names~$b$. So, for
instance, $\langle \text{Socrates}, \text{Plato}\rangle$ satisfies
$\atom{R}{x,y}$ since $\atom{R}{c,d}$ is true in the interpretation:
\begin{ekey}
	\item[\text{domain}] all people born before 2000 \textsc{ce}
	\item[a] Aristotle
	\item[b] Beyonc\'e
	\item[c] Socrates
	\item[d] Plato
	\item[\atom{P}{x}] \gap{x} is a philosopher
	\item[\atom{R}{x,y}] \gap{x} was born before \gap{y}
\end{ekey}
For atomic formulas, the objects, pairs of objects, etc., that satisfy them are exactly the extension of the predicate given in the interpretation. But the notion of satisfaction also applies to non-atomic formulas, e.g., the formula $\atom{P}{x} \land \atom{R}{x,b}$ is satisfied by all philosophers born before Beyonc\'e. It even applies to formulas involving quantifiers, e.g., $\atom{P}{x} \eand \lnot\exists y(\atom{P}{y} \land \atom{R}{y,x})$ is satisfied by all people who are philosophers and for whom it is true that no philosopher was born before them---in other words, it is true of the first philosopher.

By considering formulas (possibly involving quantifiers) with two free
variables, we can express relations for which we do not have dedicated
predicate symbols in our interpretation or symbolization key. Consider
the formula $\atom{R}{x,y}$. It expresses the relation `\gap{x} was
born before \gap{y}', since that is how we have specified its
extension. What happens if we switch the variables, i.e., consider
`$\atom{R}{y,x}$'?  A pair of objects \ntuple{$\text{y}, \text{x}$} in
the domain (i.e., a pair of people) satisfies $\atom{R}{y,x}$ if, and only if,
the reverse pair \ntuple{$\text{x}, \text{y}$} satisfies
$\atom{R}{x,y}$. In other words, $\atom{R}{y,x}$ expresses the
relation `\gap{x} was born \emph{after} \gap{y}'. Or suppose we add to
our interpretation a predicate for `teacher of'.
\begin{ekey}
	\item[\atom{T}{x,y}] \gap{x} was a teacher of \gap{y}
\end{ekey}
Then the formula `$\exists z(\atom{T}{z,x} \land \atom{T}{z,y})$' is
satisfied by x and~y if, and only if, some person z was a teacher of
both x and~y, i.e., it expresses `\gap{x} and \gap{y} have a teacher
in common'. Similarly, `$\forall z(\atom{T}{x,z} \eiff \atom{T}{y,
z})$' expresses `\gap{x} and \gap{y} taught the same people'.

The take-home message of these examples is that some English
predicates, such as `\gap{x} and \gap{y} have a teacher in common',
can be sometimes be expressed in an interpretation even if there is no
explicit predicate symbol available for them. If that is the case, you
can use a formula that expresses them (such as `$\exists
z(\atom{T}{z,x} \land \atom{T}{z,y})$') to symbolize English sentences
involving the predicate.

\practiceproblems

\problempart
Given the following interpretation:
\begin{ekey}
	\item[\text{domain}] people
	\item[\atom{W}{x}] \gap{x} is a woman (or girl)
	\item[\atom{M}{x}] \gap{x} is a man (or boy)
	\item[\atom{Y}{x,y}] \gap{x} is younger than \gap{y}
	\item[\atom{O}{x,y}] \gap{x} is an offspring of \gap{y}
	\item[d] Dana
	\item[a] Alex
\end{ekey}
express the following relations:
\begin{compactlist}
	\item \gap{x} is older than \gap{y}
	\item \gap{x} is \gap{y}'s mother
	\item \gap{x} and \gap{y} are siblings (note that you can't be your
	own sibling)
	\item \gap{x} is \gap{y}'s brother
\end{compactlist}

\problempart
Using the symbolization key from the previous
exercise, symbolize the following sentences:
\begin{compactlist}
	\item Alex and Dana are sisters.
	\item Fathers are older than their children.
	\item Alex's parents are the same age.
\end{compactlist}


\chapter{Using interpretations}
\label{sec.UsingModels}

\section{Validities and contradictions}
Suppose we want to show that `$\exists x\, \atom{A}{x,x} \eif \atom{B}{d}$' is \emph{not} a validity. This requires showing that the sentence is not true in every interpretation; i.e.,\ that it is false in some interpretation. If we can provide just one interpretation in which the sentence is false, then we will have shown that the sentence is not a validity.

In order for `$\exists x\,\atom{A}{x,x} \eif \atom{B}{d}$' to be false, the antecedent (`$\exists x\, \atom{A}{x,x}$') must be true, and the consequent (`$\atom{B}{d}$') must be false. To construct such an interpretation, we start by specifying a domain. Keeping the domain small makes it easier to specify what the predicates will be true of, so we will start with a domain that has just one member. For concreteness, let's say it is \emph{just} the city of Paris.
	\begin{ekey}
		\item[\text{domain}] Paris
	\end{ekey}
The name `$d$' must refer to something in the domain, so we have no option but:
	\begin{ekey}
		\item[d] Paris
	\end{ekey}
Recall that we want `$\exists x\, \atom{A}{x,x}$' to be true, so we want all members of the domain to be paired with themselves in the extension of `$A$'. We can just offer:
	\begin{ekey}
		\item[\atom{A}{x,y}] \gap{x} is identical with \gap{y}
	\end{ekey}
Now `$\atom{A}{d,d}$' is true, so it is surely true that `$\exists x\, \atom{A}{x,x}$'. Next, we want `$\atom{B}{d}$' to be false, so the referent of `$d$' must not be in the extension of `$B$'. We might simply offer:
	\begin{ekey}
		\item[\atom{B}{x}] \gap{x} is in Germany
	\end{ekey}
Now we have an interpretation where `$\exists x\, \atom{A}{x,x}$' is true, but where `$\atom{B}{d}$' is false. So there is an interpretation where `$\exists x\, \atom{A}{x,x} \eif \atom{B}{d}$' is false. So `$\exists x\, \atom{A}{x,x} \eif \atom{B}{d}$' is not a validity.

We can just as easily show that `$\exists x\,\atom{A}{x,x} \eif \atom{B}{d}$' is not a contradiction. We need only specify an interpretation in which `$\exists x\,\atom{A}{x,x} \eif \atom{B}{d}$' is true; i.e., an interpretation in which either `$\exists x\, \atom{A}{x,x}$' is false or `$\atom{B}{d}$' is true. Here is one:
	\begin{ekey}
		\item[\text{domain}] Paris
		\item[d] Paris
		\item[\atom{A}{x,y}] \gap{x} is identical with \gap{y}
		\item[\atom{B}{x}] \gap{x} is in France
	\end{ekey}
This shows that there is an interpretation where `$\exists x\,\atom{A}{x,x} \eif \atom{B}{d}$' is true. So `$\exists x\, \atom{A}{x,x} \eif \atom{B}{d}$' is not a contradiction.
	\factoidbox{
		\begin{factoiditems}
			\item To show that $\metav{A}$ is not a validity, it suffices to
			find an interpretation where $\metav{A}$ is false.

			\item To show that $\metav{A}$ is not a contradiction, it
			suffices to find an interpretation where $\metav{A}$ is true.
		\end{factoiditems}
	}

\section{Logical equivalence}
Suppose we want to show that `$\forall x\, \atom{S}{x}$' and `$\exists x\, \atom{S}{x}$' are not logically equivalent. We need to construct an interpretation in which the two sentences have different truth values; we want one of them to be true and the other to be false. We start by specifying a domain. Again, we make the domain small so that we can specify extensions easily. In this case, we will need at least two objects. (If we chose a domain with only one member, the two sentences would end up with the same truth value. In order to see why, try constructing some partial interpretations with one-member domains.) For concreteness, let's take:
	\begin{ekey}
		\item[\text{domain}] Ornette Coleman, Miles Davis
	\end{ekey}
We can make `$\exists x\, \atom{S}{x}$' true by including something in the extension of `$S$', and we can make `$\forall x\, \atom{S}{x}$' false by leaving something out of the extension of `$S$'. For concreteness, let's say:
	\begin{ekey}
		\item[\atom{S}{x}] \gap{x} plays saxophone
	\end{ekey}
Now `$\exists x\, \atom{S}{x}$' is true, because `$\atom{S}{x}$' is true of Ornette Coleman. Slightly more precisely, extend our interpretation by allowing `$c$' to name Ornette Coleman.  `$\atom{S}{c}$' is true in this extended interpretation, so `$\exists x\, \atom{S}{x}$' was true in the original interpretation. Similarly, `$\forall x\, \atom{S}{x}$' is false, because `$\atom{S}{x}$' is false of Miles Davis. Slightly more precisely, extend our interpretation by allowing `$d$' to name Miles Davis, and `$\atom{S}{d}$' is false in this extended interpretation, so `$\forall x\, \atom{S}{x}$' was false in the original interpretation. We have provided a counter-interpretation to the claim that `$\forall x\, \atom{S}{x}$' and `$\exists x\, \atom{S}{x}$' are logically equivalent.
	\factoidbox{
		To show that $\metav{A}$ and $\metav{B}$ are not logically equivalent, it suffices to find an interpretation where one is true and the other is false.
	}

\section{Validity, entailment and satisfiability}
To test for validity, entailment, or satisfiability, we typically need to produce interpretations that determine the truth value of several sentences simultaneously.

Consider the following argument in FOL:
$$\exists x(\atom{G}{x} \eif \atom{G}{a}) \therefore \exists x\, \atom{G}{x} \eif \atom{G}{a}$$
To show that this is invalid, we must make the premise true and the conclusion false. The conclusion is a conditional, so to make it false, the antecedent must be true and the consequent must be false. Clearly, our domain must contain two objects. Let's try:
	\begin{ekey}
		\item[\text{domain}] Karl Marx, Ludwig von Mises
		\item[\atom{G}{x}] \gap{x} hated communism
		\item[a] Karl Marx
	\end{ekey}
Given that Marx wrote \emph{The Communist Manifesto}, `$\atom{G}{a}$' is plainly false in this interpretation. But von Mises famously hated communism, so `$\exists x\, \atom{G}{x}$' is true in this interpretation. Hence `$\exists x\, \atom{G}{x} \eif \atom{G}{a}$' is false, as required.

Does this interpretation make the premise true? Yes it does! Note that `$\atom{G}{a} \eif \atom{G}{a}$' is true. (Indeed, it is a validity.) But then certainly `$\exists x (\atom{G}{x} \eif \atom{G}{a})$' is true, so the premise is true, and the conclusion is false, in this interpretation. The argument is therefore invalid.

In passing, note that we have also shown that `$\exists x(\atom{G}{x} \eif \atom{G}{a})$' does \emph{not} entail `$\exists x\, \atom{G}{x} \eif \atom{G}{a}$', i.e., that $\exists x (\atom{G}{x} \eif \atom{G}{a}) \nentails \exists x\, \atom{G}{x} \eif \atom{G}{a}$. Equally, we have shown that the sentences `$\exists x (\atom{G}{x} \eif \atom{G}{a})$' and `$\enot (\exists x\, \atom{G}{x} \eif \atom{G}{a})$' are jointly satisfiable.

Let's consider a second example:
$$\forall x \exists y\, \atom{L}{x,y} \therefore \exists y \forall x\, \atom{L}{x,y}$$
Again, we want to show that this is invalid. To do this, we must make the premises true and the conclusion false. Here is a suggestion:
\begin{ekey}
	\item[\text{domain}] Canadian citizens currently in a domestic partnership with another Canadian citizen
	\item[\atom{L}{x,y}] \gap{x} is in a domestic partnership with \gap{y}
\end{ekey}
The premise is clearly true on this interpretation. Anyone in the domain is a Canadian citizen in a domestic partnership with some other Canadian citizen. That other citizen will also, then, be in the domain. So for everyone in the domain, there will be someone (else) in the domain with whom they are in a domestic partnership. Hence `$\forall x \exists y\, \atom{L}{x,y}$' is true. However, the conclusion is clearly false, for that would require that there is some single person who is in a domestic partnership with everyone in the domain, and there is no such person, so the argument is invalid. We observe immediately that the sentences `$\forall x \exists y\, \atom{L}{x,y}$' and `$\enot\exists y \forall x\, \atom{L}{x,y}$' are jointly satisfiable and that `$\forall x \exists y\, \atom{L}{x,y}$' does not entail `$\exists y \forall x\, \atom{L}{x,y}$'.

For our third example, we'll mix things up a bit. In \cref{s:Interpretations}, we described how we can present some interpretations using diagrams. For example:
\begin{center}
	\begin{tikzpicture}[alt={The numbers 1, 2, and 3 are connected by arrows. There are arrows from 1 to 2, 3, and to itself, and from 2 to 3. There is no arrow leaving 3.}]
	\node (atom1) at (0,2) {1};
	\node (atom2) at (2,2) {2};
	\node (atom3) at (2,0) {3};
	\draw[->, thick] (atom1)--(atom2);
	\draw[->, thick] (atom1)--(atom3);
	\draw[->, thick] (atom1)+(-0.15,0.15) arc (-330:-30:.3); 
	\draw[->, thick] (atom2) -- (atom3);
\end{tikzpicture}
\bmlDescription{The numbers 1, 2, and 3 are connected by arrows. There are arrows from 1 to 2, 3, and to itself, and from 2 to 3. There is no arrow leaving 3.}
\end{center}
Using the conventions employed in \cref{s:Interpretations}, the domain of this interpretation is the first three positive whole numbers, and `$\atom{R}{x,y}$' is true of x and y just in case there is an arrow from x to y in our diagram. Here are some sentences that the interpretation makes true:
\begin{itemize}
	\item `$\forall x \exists y\, \atom{R}{y,x}$' 
	\item `$\exists x \forall y\, \atom{R}{x,y}$' \hfill (witness for
	$x$: $1$)
	\item `$\exists x \forall y (\atom{R}{y,x} \eiff x = y)$' \hfill (witness for
	$x$: $1$)
	\item `$\exists x \exists y \exists z ((\enot y = z \eand \atom{R}{x,y}) \eand \atom{R}{z,x})$' \hfill (witness for
	$x$: $2$)
	\item `$\exists x \forall y\, \enot \atom{R}{x,y}$' \hfill (witness for
	$x$: $3$)
	\item `$\exists x (\exists y\, \atom{R}{y,x} \eand \enot \exists y\, \atom{R}{x,y})$' \hfill (witness for
	$x$: $3$)
\end{itemize}
This immediately shows that all of the preceding six sentences are jointly satisfiable. We can use this observation to generate \emph{invalid} arguments, e.g.:
\begin{align*}
	\forall x \exists y\, \atom{R}{y,x}, \exists x \forall y\, \atom{R}{x,y}  &\therefore  \forall x \exists y\, \atom{R}{x,y}\\
	\exists x \forall y\, \atom{R}{x,y}, \exists x \forall y \enot \atom{R}{x,y} & \therefore \enot \exists x \exists y \exists z (\enot y = z \eand (\atom{R}{x,y} \eand \atom{R}{z,x}))
\end{align*}
and many more besides.

\factoidbox{
	If some interpretation makes all of $\metav{A}_1, \metav{A}_2, \ldots, \metav{A}_n$  true and $\metav{C}$ is false, then:
	\begin{compactlist}
		\item$\metav{A}_1, \metav{A}_2, \ldots, \metav{A}_n \therefore \metav{C}$ is \emph{invalid}; and
	\item $\metav{A}_1, \metav{A}_2, \ldots, \metav{A}_n \nentails \metav{C}$; and
	\item 
	And $\metav{A}_1, \metav{A}_2, \ldots, \metav{A}_n, \enot \metav{C}$ are jointly satisfiable.
\end{compactlist}}
An interpretation which refutes a claim---to logical truth, say, or to entailment---is called a \emph{counter-interpretation}, or a \emph{counter-model}.

We'll close this section, though, with a caution about the relationship between (in)validity and (non)entailment. Recall FOL's limitations: it is an extensional language; it ignores issues of vagueness; and it cannot handle cases of validity for `special reasons'. To take one illustration of these issues, consider this natural-language argument: 
\begin{earg}
	\item[] Every fox is cute.
	\item[\texttherefore] All vixens are cute.
\end{earg}
This is conceptually valid: necessarily every vixen is a fox, so it is impossible for the premise to be true and the conclusion false. Now, we might sensibly symbolize the argument as follows:
$$\forall x(\atom{F}{x} \eif \atom{C}{x}) \therefore \forall x(\atom{V}{x} \eif  \atom{C}{x})$$
However, it is easy to find counter-models which show that $\forall x(\atom{F}{x} \eif \atom{C}{x}) \nentails \forall x(\atom{V}{x} \eif  \atom{C}{x})$. (\emph{Exercise}: find one.) So, it would be \emph{wrong} to infer that the English argument is \emph{invalid}, just because there is a counter-model to the relevant \emph{FOL{}-entailment}.

The general moral is this. If you want to infer from the absence of an entailment in FOL to the invalidity of some English argument, then you need to argue that nothing important is lost in the way you have symbolized the English argument.

\practiceproblems

\solutions
\problempart
\label{pr.Contingent}
Show that each of the following is neither a validity nor a contradiction:
\begin{compactlist}
\item \leftsolutions\ $\atom{D}{a}  \eand \atom{D}{b}$
\item \leftsolutions\ $\exists x\, \atom{T}{x,h}$
\item \leftsolutions\ $\atom{P}{m}  \eand \enot\forall x\, \atom{P}{x}$
\item $\forall z \atom{J}{z} \eiff \exists y\, \atom{J}{y}$
\item $\forall x (\atom{W}{x,m,n} \eor \exists y\,\atom{L}{x,y})$
\item $\exists x (\atom{G}{x} \eif \forall y\, \atom{M}{y})$
\item $\exists x (x = h \eand x = i)$
\end{compactlist}

\solutions
\problempart
\label{pr.NotEquiv}
Show that the following pairs of sentences are not logically equivalent.
\begin{compactlist}
\item $\atom{J}{a} $,  $\atom{K}{a}$
\item $\exists x\, \atom{J}{x}$,  $\atom{J}{m}$
\item $\forall x\, \atom{R}{x,x}$, $\exists x\, \atom{R}{x,x}$
\item $\exists x\, \atom{P}{x} \eif \atom{Q}{c}$, $\exists x (\atom{P}{x} \eif \atom{Q}{c})$
\item $\forall x(\atom{P}{x} \eif \enot \atom{Q}{x})$, $\exists x(\atom{P}{x} \eand \enot \atom{Q}{x})$
\item $\exists x(\atom{P}{x} \eand \atom{Q}{x})$, $\exists x(\atom{P}{x} \eif \atom{Q}{x})$
\item $\forall x(\atom{P}{x}\eif \atom{Q}{x})$, $\forall x(\atom{P}{x} \eand \atom{Q}{x})$
\item $\forall x\exists y\, \atom{R}{x,y}$, $\exists x\forall y\, \atom{R}{x,y}$
\item $\forall x\exists y\, \atom{R}{x,y}$, $\forall x\exists y\, \atom{R}{y,x}$
\end{compactlist}



\problempart
Show that the following sentences are jointly satisfiable:
\begin{compactlist}
\item  $\atom{M}{a}, \enot \atom{N}{a}, \atom{P}{a}, \enot \atom{Q}{a}$
\item $\atom{L}{e,e}, \atom{L}{e,g}, \enot \atom{L}{g,e}, \enot \atom{L}{g,g}$
\item $\enot (\atom{M}{a} \eand \exists x\, \atom{A}{x}), \atom{M}{a} \eor \atom{F}{a}, \forall x(\atom{F}{x} \eif \atom{A}{x})$
\item $\atom{M}{a} \eor \atom{M}{b}, \atom{M}{a} \eif \forall x \enot \atom{M}{x}$
\item $\forall y\, \atom{G}{y}, \forall x (\atom{G}{x} \eif \atom{H}{x}), \exists y \enot \atom{I}{y}$
\item $\exists x(\atom{B}{x} \eor \atom{A}{x}), \forall x \enot \atom{C}{x}, \forall x\bigl[(\atom{A}{x} \eand \atom{B}{x}) \eif \atom{C}{x}\bigr]$
\item $\exists x\, \atom{X}{x}, \exists x\, \atom{Y}{x}, \forall x(\atom{X}{x} \eiff \enot \atom{Y}{x})$
\item $\forall x(\atom{P}{x} \eor \atom{Q}{x}), \exists x\enot(\atom{Q}{x} \eand \atom{P}{x})$
\item $\exists z(\atom{N}{z} \eand \atom{O}{z,z}), \forall x\forall y(\atom{O}{x,y} \eif \atom{O}{y,x})$
\item $\enot \exists x \forall y\, \atom{R}{x,y}, \forall x \exists y\, \atom{R}{x,y}$
\item $\enot \atom{R}{a,a}$, $\forall x (x=a \eor \atom{R}{x,a})$
\item $\forall x\forall y\forall z[(x=y \eor y=z )\eor x=z]$, $\exists x\exists y\ \enot x= y$
\item $\exists x\exists y((\atom{Z}{x} \eand \atom{Z}{y} )\eand x=y)$, $\enot \atom{Z}{d}$, $d=e$
\end{compactlist}

\problempart
Show that the following arguments are invalid:
\begin{compactlist}
\item $\forall x(\atom{A}{x} \eif \atom{B}{x}) \therefore \exists x\, \atom{B}{x}$
\item $\forall x(\atom{R}{x} \eif \atom{D}{x}), \forall x(\atom{R}{x} \eif \atom{F}{x}) \therefore \exists x(\atom{D}{x} \eand \atom{F}{x})$
\item $\exists x(\atom{P}{x}\eif \atom{Q}{x}) \therefore \exists x\, \atom{P}{x}$
\item $\atom{N}{a} \eand \atom{N}{b} \eand \atom{N}{c} \therefore \forall x\, \atom{N}{x}$
\item $\atom{R}{d,e}, \exists x\, \atom{R}{x,d} \therefore \atom{R}{e,d}$
\item $\exists x(\atom{E}{x} \eand \atom{F}{x}), \exists x\, \atom{F}{x} \eif \exists x\, \atom{G}{x} \therefore \exists x(\atom{E}{x} \eand \atom{G}{x})$
\item $\forall x\, \atom{O}{x,c}, \forall x\, \atom{O}{c,x} \therefore \forall x\, \atom{O}{x,x}$
\item $\exists x(\atom{J}{x} \eand \atom{K}{x}), \exists x \enot \atom{K}{x}, \exists x \enot \atom{J}{x} \therefore \exists x(\enot \atom{J}{x} \eand \enot \atom{K}{x})$
\item $\atom{L}{a,b} \eif \forall x\, \atom{L}{x,b}, \exists x\, \atom{L}{x,b} \therefore \atom{L}{b,b}$
\item $\forall x(\atom{D}{x} \eif \exists y\, \atom{T}{y,x}) \therefore \exists y \exists z\ \enot y= z$
\end{compactlist}

\chapter{Reasoning about interpretations}

\section{Validities and contradictions}
We can show that a sentence is \emph{not} a validity just by providing one carefully specified interpretation: an interpretation in which the sentence is false. To show that something \emph{is} a validity, on the other hand, it would not be enough to construct ten, one hundred, or even a thousand interpretations in which the sentence is true. A sentence is only a validity if it is true in \emph{every} interpretation, and there are infinitely many interpretations. We need to reason about all of them, and we cannot do this by dealing with them one by one!

Sometimes, we can reason about all interpretations fairly easily. For example, we can offer a relatively simple argument that `$\atom{R}{a,a}\eor\enot \atom{R}{a,a}$' is a validity:
	\begin{quote}
		\label{allmodels1}
		Any relevant interpretation will give `$\atom{R}{a,a}$' a truth value. If `$\atom{R}{a,a}$' is true in an interpretation, then `$\atom{R}{a,a} \eor \enot\atom{R}{a,a}$' is true in that interpretation. If `$\atom{R}{a,a}$' is false in an interpretation, then $\enot\atom{R}{a,a}$ is true, and so `$\atom{R}{a,a} \eor\enot \atom{R}{a,a}$' is true in that interpretation. These are the only alternatives. So `$\atom{R}{a,a} \eor\enot \atom{R}{a,a}$' is true in every interpretation. Therefore, it is a validity.
	\end{quote}
This argument is valid, of course, and its conclusion is true. However, it is not an argument in FOL. Rather, it is an argument in English \emph{about} FOL: it is an argument in the metalanguage.

Note another feature of the argument. Since the sentence in question contained no quantifiers, we did not need to think about how to interpret `$a$' and `$R$'; the point was just that, however we interpreted them, `$\atom{R}{a,a}$' would have some truth value or other. (We could ultimately have given the same argument concerning TFL sentences.)

Let's have another example. The sentence `$\forall x(\atom{R}{x,x}\eor\enot \atom{R}{x,x})$' should obviously be a validity. However, saying \emph{precisely} why is quite tricky. We cannot say that `$\atom{R}{x,x} \eor\enot \atom{R}{x,x}$' is true in every interpretation, since `$\atom{R}{x,x} \eor\enot \atom{R}{x,x}$' is not even a \emph{sentence} of FOL (remember that `$x$' is a variable, not a name). Instead, we should say something like this:
	\begin{quote}
		Consider some arbitrary interpretation. `$\forall x(\atom{R}{x,x}\eor \enot\atom{R}{x,x})$' is true in our interpretation \ifeff{} `$\atom{R}{x,x}\eor\enot\atom{R}{x,x}$' is satisfied by every object of its domain. Consider some arbitrary member of the domain, which, for convenience, we will call Fred. Either Fred satisfies `$\atom{R}{x,x}$' or it does not. If Fred satisfies `$\atom{R}{x,x}$', then Fred also satisfies `$\atom{R}{x,x} \eor \enot \atom{R}{x,x}$'. If Fred does not satisfy `$\atom{R}{x,x}$', it \emph{does} satisfy `$\enot\atom{R}{x,x}$' and so also `$\atom{R}{x,x} \eor\enot \atom{R}{x,x}$'.\footnote{We use here the fact that the truth conditions for connectives also apply to satisfaction: $a$ satisfies $\metav{A}(\metav{x}) \lor \metav{B}(\metav{x})$ \ifeff{} $a$ satisfies $\metav{A}(\metav{x})$ or $\metav{B}(\metav{x})$, etc.} So either way, Fred satisfies `$\atom{R}{x,x} \eor\enot \atom{R}{x,x}$'. Since there was nothing special about Fred---we might have chosen any object---we see that every object in the domain satisfies `$\atom{R}{x,x} \eor\enot \atom{R}{x,x}$'. So `$\forall x (\atom{R}{x,x} \eor\enot \atom{R}{x,x})$' is true in our interpretation. But we chose our interpretation arbitrarily, so `$\forall x (\atom{R}{x,x} \eor\enot \atom{R}{x,x})$' is true in every interpretation. It is therefore a validity.
	\end{quote}
This is quite long-winded, but, as things stand, there is no alternative. In order to show that a sentence is a validity, we must reason about \emph{all} interpretations.

\section{Other cases}
Similar points hold of other cases too. Thus, we must reason about all interpretations if we want to show:
	\begin{itemize}
		\item that a sentence is a contradiction (this requires that it is false in \emph{every} interpretation);
		\item that two sentences are logically equivalent (this requires that they have the same truth value in \emph{every} interpretation);
		\item that some sentences are jointly unsatisfiable (this requires that there is no interpretation in which all of those sentences are true together, i.e., that, in \emph{every} interpretation, at  least one of those sentences is false);
		\item that an argument is valid (this requires that the conclusion is true in \emph{every} interpretation where the premises are true);
		\item that some sentences entail another sentence.
	\end{itemize}
The problem is that, with the tools available to you so far, reasoning about all interpretations is a serious challenge! For a final example, here is a perfectly obvious entailment:
	$$\forall x(\atom{H}{x} \eand \atom{J}{x}) \entails \forall x\, \atom{H}{x}$$
After all, if everything is both $H$ and $J$, then everything is~$H$.
But we can only establish the entailment by considering what must be
true in every interpretation in which the premise $\forall
x(\atom{H}{x} \eand \atom{J}{x})$ is true. To show this, we would have
to reason as follows:
	\begin{quote}
		Consider an arbitrary interpretation in which `$\forall x(\atom{H}{x} \eand \atom{J}{x})$' is true. It follows that `$\atom{H}{x} \eand \atom{J}{x}$' is satisfied by every object in this interpretation. `$\atom{H}{x}$' will, then, also be satisfied by every object.\footnote{Here again we make use of the fact that any object that satisfies $\metav{A}(\metav{x}) \land \metav{B}(\metav{x})$ must satisfy both $\metav{A}(\metav{x})$ and $\metav{B}(\metav{x})$.} So it must be that `$\forall x\, \atom{H}{x}$' is true in the  interpretation. We've assumed nothing about the interpretation except that it was one in which `$\forall x(\atom{H}{x} \eand \atom{J}{x})$' is true. So any interpretation in which `$\forall x(\atom{H}{x} \eand \atom{J}{x})$' is true is one in which `$\forall x\, \atom{H}{x}$' is true.
\end{quote}
Even for a simple entailment like this one, the reasoning is somewhat complicated. For more complicated entailments, the reasoning can be extremely torturous.

The following table summarizes whether a single interpretation or counter-interpretation suffices, or whether we must reason about all interpretations.

\begin{center}\small
\begin{tabular}{l l l}
%\cline{2-3}
 & \textbf{Yes} & \textbf{No}\\
 \hline
%\cline{2-3}
validity? & all interpretations & one counter-interpretation\\
contradiction? &  all interpretations  & one counter-interpretation\\
equivalent? & all interpretations & one counter-interpretation\\
satisfiable? & one interpretation & all interpretations\\
valid? & all interpretations & one counter-interpretation\\
entailment? & all interpretations & one counter-interpretation\\
\end{tabular}
\end{center}
\label{table.ModelOrArgument}

You might want to compare this table with the table at the end of \cref{s:PartialTruthTable}. The key difference resides in the fact that TFL concerns truth tables, whereas FOL concerns interpretations. This difference is deeply important, since each truth-table only ever has finitely many lines, so that a complete truth table is a relatively tractable object. By contrast, there are infinitely many interpretations for any given sentence(s), so that reasoning about all interpretations can be a deeply tricky business.

\chapter{Properties of relations}\label{ch:PropRelations}

Symbolization keys allow us to assign two-place predicate symbols like
`$\atom{R}{x,y}$' to English predicates with two gaps, e.g., `\gap{x}
is older than \gap{y}'. This allows us to symbolize arguments such as:
\begin{earg}
	\item Rudolf is older than Kurt.
	\item Kurt is older than Julia.
	\item[\texttherefore] Rudolf is older than Julia.
\end{earg}
This argument is valid, but its symbolization in FOL,
\begin{earg}
	\item $\atom{R}{r,k}$
	\item $\atom{R}{k,j}$
	\item[\texttherefore] $\atom{R}{r,j}$
\end{earg}
is not. That's because the validity of the argument depends on a
property of the relation `$x$ is older than $y$', namely that if some
person $x$ is older than a person~$y$, and $y$ is also older than~$z$,
then $x$ must be older than $z$. This property of `older than' is
called \define{transitivity}. We can symbolize this property
\emph{itself} in FOL as:
\[
	\forall x\forall y\forall z((\atom{R}{x,y} \eand \atom{R}{y,z}) \eif
	\atom{R}{x,z})
\]
Whenever the validity of an argument \emph{only} depends on a property
of a relation involved, and this property can be symbolized in FOL, we
can add this symbolization to the argument and obtain an argument that
is in fact valid in FOL:
\begin{earg}
	\item $\forall x\forall y\forall z((\atom{R}{x,y} \eand \atom{R}{y,z}) \eif
	\atom{R}{x,z})$
	\item $\atom{R}{r,k}$
	\item $\atom{R}{k,j}$
	\item[\texttherefore] $\atom{R}{r,j}$
\end{earg}

Properties of relations such as transitivity are important concepts
especially in applications of logic in the sciences. For instance,
order relations between numbers such as $<$ and $\le$ (and also $>$
and $\ge$) are transitive. The identity relation~$=$ is also
transitive.

There are other properties of relations that are important enough to
have names, and often show up in applications. Here are some:

\factoidbox{
	\begin{factoiditems}
		\item A relation $R$ is \define{transitive} \ifeff{} it is the case
		that whenever $\atom{R}{x,y}$ and $\atom{R}{y,z}$ then also~$\atom{R}{x,z}$.
		\item A relation $R$ is \define{reflexive} \ifeff{} for any~$x$,
		$\atom{R}{x,x}$ holds.
		\item A relation $R$ is \define{symmetric} \ifeff{} whenever
		$\atom{R}{x,y}$ holds, so does $\atom{R}{y,x}$.
		\item A relation $R$ is \define{anti-symmetric} \ifeff{} for no two
		\emph{different} $x$ and $y$, $\atom{R}{x,y}$ and $\atom{R}{y,x}$
		both hold.
	\end{factoiditems}}

\newglossaryentry{transitive}
{
  name=transitivity,
  text = transitive,
  description={The property a relation~$R$ has \ifeff{} it is the case
  that whenever $\atom{R}{x,y}$ and $\atom{R}{y,z}$ then also~$\atom{R}{x,z}$}
}
\newglossaryentry{reflexive}
{
  name=reflexivity,
  text = reflexive,
  description={The property a relation~$R$ has \ifeff{} for any~$x$, $\atom{R}{x,x}$ holds}
}
\newglossaryentry{symmetric}
{
  name=symmetry,
  text = symmetric,
  description={The property a relation~$R$ has \ifeff{} it is the case
  that whenever $\atom{R}{x,y}$ holds, so does $\atom{R}{y,x}$}
}
\newglossaryentry{anti-symmetric}
{
  name=anti-symmetry,
  text = refleanti-symmetricxive,
  description={The property a relation~$R$ has \ifeff{} for no two
	\emph{different} $x$ and $y$, $\atom{R}{x,y}$ and $\atom{R}{y,x}$
	both hold}
}

We have already seen that transitivity can be symbolized in FOL. The
others can, too:
\begin{itemize}
	\item $R$ is reflexive: $\forall x\,\atom{R}{x,x}$
	\item $R$ is symmetric: \[\forall x\forall y(\atom{R}{x,y} \eif \atom{R}{y,x})\]
	\item $R$ is anti-symmetric: \[\lnot\exists x\exists y((\atom{R}{x,y}
	\eand \atom{R}{y,x}) \eand \enot x=y)\] or, equivalently, \[\forall x\forall y((\atom{R}{x,y}
	\eand \atom{R}{y,x}) \eif x=y)\]
\end{itemize}

Relations expressing an equivalence in some respect are reflexive,
e.g., `is the same age as', `is as tall as', and the most stringent
equivalence of them all, identity~$=$. They are also symmetric: e.g.,
whenever $x$ is as tall as~$y$, then $y$ is as tall as~$x$. In fact,
relations that are reflexive, symmetric, and transitive are called
\define{equivalence relations}.

Equivalences aren't the only symmetric relations. For instance, `is a
sibling of' is symmetric, but it is not reflexive (as no one is their own
sibling).

Relations that are reflexive, transitive, and anti-symmetric are
called \define{partial orders}. For instance, $\le$ and $\ge$ are
partial orders. The relation `is no older than', by contrast, is
reflexive and transitive, but not anti-symmetric. (Two different
people can be of the same age, and so neither is older than the
other.)

\practiceproblems

\problempart
Give examples of relations with the following properties:
\begin{compactlist}
	\item Reflexive but not symmetric
	\item Symmetric but not transitive
	\item Transitive, symmetric, but not reflexive
\end{compactlist}

\problempart
Show that a relation can be both symmetric and anti-symmetric by
giving an interpretation that makes both of the following sentences
true:
\begin{align*}
	& \forall x\forall y(\atom{R}{x,y} \eif \atom{R}{y,x})\\
	& \forall x\forall y((\atom{R}{x,y} \eand \atom{R}{y,x}) \eif x=y)
\end{align*}

